\chapter{Who Bears the Transition}
\label{ch:managing-moment}

\epigraph{A machine never decides who keeps the hours it saves. People decide that, and where they decline to decide, whoever moves first decides it for them.}{}

This book has spent its first part diagnosing a coordination collapse and describing substrates meant to survive it. It has not yet said why the moment has to be \emph{managed} rather than merely endured. The case is easier to make from two hundred years' distance than from inside the noise of 2026, so this chapter starts in a Yorkshire valley in April 1812, where a group of men who are now a byword for stupidity were asking, in writing and under law, for exactly the kind of managed transition this book argues for. They were refused one. What happened next is the strongest available argument that refusing one is the expensive option.

\section{Before the Frames Broke}

Before a single frame broke in the West Riding, the croppers there had spent years asking for a managed transition on paper. Between 1802 and 1806 they petitioned Parliament and their employers under the old apprenticeship and gig-mill statutes, and their proposals read as transition policy. They asked that the new shearing machinery be brought in gradually; that men displaced by it be found alternative work; and that a levy of sixpence on every yard of machine-dressed cloth be set aside to support the unemployed while they looked for it \parencite{Thompson1963,Randall1991}. This was a constitutional campaign, conducted through the channels a subject of the Crown was supposed to use, and it was pragmatic to the point of modesty. It asked for a share of the machine's yield and a pace the trade could absorb.

The campaign exhausted itself against a closing door. The croppers had carried their case as far as a parliamentary committee, whose 1806 report on the state of the woollen manufacture heard them out and settled nothing in their favour. Parliament then repealed the protective statutes in 1809, removing the legal ground the petitions had stood on. The Combination Acts of 1799 and 1800 had already made it a crime for workmen to associate to improve their conditions, and an older paternalist protection for the trade had been suspended and left to lapse. A Commons inquiry into the framework knitters turned, in the hosts' phrase, into an inquisition, reading the workers' organising as the first stir of revolution and advising Parliament to act against it \parencite{LynskeyDunt2026Luddites}. By 1811 the croppers held a grievance, no statute, no union, and no petition left to file. The machine-breaking began in Nottingham that year and spread into Yorkshire and Lancashire through 1812, in that order: the argument was lost by constitutional means first, and continued by sledgehammer only afterwards.

The men who took up the sledgehammer were not against machinery as such. The letters and oaths of the 1811--13 Luddites, collected by \textcite{Binfield2004}, set out demands about who ran the machines and to what standard: an end to the employment of ``colts'', men who had not served the seven-year apprenticeship; an end to fraudulent ``cut-up'' goods run off cheaply to undercut price and quality; and restoration of the customary ``Standard Price'' for piece-work. A letter signed ``General Ludd'' to a Yorkshire frame-maker objected to ``all Machinery hurtful to Commonality'', to machinery that harmed the common good by putting out shoddy work with unapprenticed hands at cut wages. Read the demands and the caricature falls apart on contact. These were people trying to negotiate the terms of an automation rather than to abolish it, and the terms they named were apprenticeship, quality, and a customary price, the ordinary furniture of a trade defending itself.

Rawfolds Mill is where the negotiation ended in gunfire. On the night of 11--12 April 1812 roughly a hundred and fifty men gathered at the Dumb Steeple and marched on William Cartwright's mill near Cleckheaton, which held shearing frames and, this time, a defence. Cartwright had been warned and had garrisoned the place with a handful of workmen and a small detachment of Cumberland Militia. The attackers hammered the doors with a great sledgehammer the croppers had named Enoch and fired through gaps in the wall; the defenders fired back; after some twenty minutes George Mellor gave the order to withdraw. Two Luddites, Samuel Hartley and John Booth, were left mortally wounded and died soon after. None of the defenders was killed \parencite{Thompson1963}.

The most-repeated detail of that night is a deathbed exchange in which the dying Booth lures a clergyman-magistrate close with the promise of a secret, then whispers only that he too can keep one, and expires. It is a wonderful line and it is probably not true. The podcast that carries it flags its own sources as including a ``non-fiction novel'' the authors themselves sold as unreliable, and the hosts describe their method as a maximal reading of the evidence, believing the best stories on offer \parencite{LynskeyDunt2026Luddites}. A book that grades its own claims elsewhere should say the same here: the frames broke and men died, and the good lines around those facts are decoration. The historiographic honesty is the point of raising them at all.

\section{A Word Written by the Winners}

Two centuries later ``Luddite'' means a person too foolish or frightened to understand technology. That definition is a victory for one side of a nineteenth-century labour dispute, and it is worth asking who it serves. It serves anyone who would like the automation question settled without discussing who captures the gains, because a Luddite, so defined, has no argument to answer. The word does the foreclosing before the debate begins.

The insult is younger than the events. As a general term for a technophobe it surfaces only around 1950 and takes off quickly, though the Oxford English Dictionary did not admit that sense until 1989; for something like a hundred and thirty years after the trials the word lay dormant, then revived as a slur \parencite{LynskeyDunt2026Luddites}. Its afterlife is the tell. Dunt puts the through-line bluntly: the propaganda used against the Luddites at the time has reached the present ``without any distillation whatsoever''. The winners wrote the caption and it stuck.

Historians have spent seventy years peeling it off. \textcite{Hobsbawm1952} reopened the question in \emph{Past \& Present}, arguing that machine-breaking was a rational technique of trade unionism, a substitute for the strike in trades that had no stable union to call one. \textcite{Thompson1963} went further, setting out to rescue the stockinger and the cropper from what he called the enormous condescension of posterity and reading Luddism as a disciplined movement rather than a spasm. The recovery continues in the present academic literature: \textcite{CharitsisLaamanenLehtiniemi2024} reclaim ``algorithmic Luddism'' as a framework for data-capitalism labour politics, built on refusal, resistance, and re-imagining, and \textcite{Merchant2023} traces a direct line from the Regency mill to the platform economy for a general readership.

Recovery carries its own hazard, and the honest sources name it. The same podcast that corrects the caricature warns against replacing it with a romance in which every frame-breaker was a saint. Real movements hold organisers, opportunists, and the merely confused at the same time, and a chapter that leaned on the Luddites as pure heroes would be committing the mirror-image error to the one it complains about. The reason to get the reading right is narrow and practical. The caricature is not merely unfair; it deletes the only policy lesson the episode carries, which is what a state does when a technology makes a trade redundant and the people in that trade ask for terms. Even Ricardo, no friend of the machine-breakers, came within a decade to concede in the third edition of his \emph{Principles} that machinery could indeed injure the interests of the labouring class \parencite{Ricardo1821}, which is to say the Luddites had read their own economy more accurately than the men who hanged them. The lesson survives only if the caricature is set aside long enough to hear the question they asked.

\section{Collective Bargaining by Riot}

Hobsbawm gave that lesson its name. He called the Luddite method ``collective bargaining by riot'': a way of exerting pressure on employers that stood in for the collective bargaining the law forbade \parencite{Hobsbawm1952}. The phrase is exact because it identifies the violence as a bargaining channel of last resort, opened only after the others were shut. By 1812 they were all shut. Association was criminal under the Combination Acts. An older paternalist protection, a statutory minimum for weavers, had been suspended and then repealed as laissez-faire hardened into orthodoxy. Petitioning had failed and been read as sedition. Dunt draws the moral from Gravener Henson's warning to the manufacturers of the day and calls it that outright: ``if you thwart hope of reform then there will be violence'' \parencite{LynskeyDunt2026Luddites}. The violence was the predicted output of a system with no legitimate input left.

Faced with that, the state chose suppression over adjustment, and it chose knowingly. In February 1812 the Frame Breaking Bill made the destruction of stocking frames a capital felony \parencite{FrameBreakingAct1812}. Lord Byron rose in the Lords on 27 February to oppose it in his maiden speech, and the passage that survives is a direct challenge to the logic of hanging one's way out of an economic problem: ``Can you commit a whole county to their own prisons? Will you erect a gibbet in every field and hang up men like scarecrows?'' \parencite{Byron1812}. The bill became law in March regardless. The alternative on the table, a transition managed by statute, had already been refused; the alternative chosen was the rope.

\begin{keyinsight}[title={Refused, Not Absent by Accident}]
The adjustment mechanisms the Luddites needed were named, costed, and proposed years before the first frame broke: phased introduction, alternative employment, a ring-fenced levy on the machine's output to fund the displaced. Parliament did not fail to think of them. It considered the constitutional demand, repealed the protection the workers stood on, criminalised their association, and then made breaking a frame a hanging offence. The transition was not missing for want of imagination. It was declined.
\end{keyinsight}

The military response is a place where a careful book has to resist a good statistic. The figure often quoted is that more troops were deployed against the Luddites than Wellington commanded in the Peninsula. Roughly twelve thousand soldiers were sent into the disturbed districts, a real and large number traced to Darvall through Hobsbawm's later work \parencite{Darvall1934,Hobsbawm1964}. The comparison with Wellington, however, holds only against Wellington's smallest force, the expeditionary body that first landed in 1808; by 1812 Britain had something closer to fifty thousand men in the Peninsula, and the home troops were largely Militia and Yeomanry legally barred from serving abroad, so they were never Napoleon's to lose \parencite{Howes2017,Linch2011}. The occupation was heavy and the twelve thousand is defensible; the Wellington line is a cherry-pick and this book will not print it as fact.

What followed the occupation was proportionate to nothing. The York Special Commission sat in January 1813. Three men, Mellor among them, were hanged on the eighth for the murder of the manufacturer William Horsfall; fourteen more went to the gallows on the sixteenth for Rawfolds and related offences, seventeen executions in eight days \parencite{YorkProceedings1813,Thompson1963}. The trade the machinery displaced did not adjust so much as vanish. By the episode's account Yorkshire held around five thousand croppers in 1812 and fewer than a thousand five years on, the remnant working for a quarter of their former wages \parencite{LynskeyDunt2026Luddites}. That local collapse sat inside a national one: \textcite{Allen2009} shows British real wages stagnating across roughly the first half of the century while output per worker and the profit share climbed, an ``Engels' pause'' of fifty to seventy years in which the gains of mechanisation went one way and the costs went the other.

The suppression was also, on any accounting, the expensive way to lose. Kirkpatrick Sale reckoned the frames and shears the Luddites actually destroyed at something near a hundred thousand pounds; the military occupation, the guarding of factories, the spies and the trials cost the state a multiple of that, so the response came dearer than the machinery it defended \parencite{LynskeyDunt2026Luddites}. A managed transition, a levy and a phasing, would have been cheaper than the garrison. The reform that eventually reached these workers' successors was constitutional, the repeal of the Combination Acts in 1824 and the Reform Act of 1832, and it arrived a generation too late for the men hanged at York. The state had meanwhile, in the hosts' summary, forged a lasting alliance with capital and learned to meet the next rebellion with soldiers, spies, and the scaffold \parencite{LynskeyDunt2026Luddites}.

\section{Reverse-Centaurs and Funny Math}

The present rhymes with this, and the sharpest guide to the rhyme is Cory Doctorow, whose 2026 book supplies a vocabulary for the same fight in its current costume \parencite{Doctorow2026ReverseCentaur}. His central image is the reverse-centaur, a term he coined years earlier and has since made load-bearing \parencite{Doctorow2021ReverseCentaur}.

\begin{definitionbox}[title={Centaur and Reverse-Centaur}]
A \textbf{centaur} is a human assisted by a machine: the person chooses when the tool runs, keeps their judgement in the loop, and is served by the automation. A \textbf{reverse-centaur} is the inversion. In Doctorow's words, ``a human who is conscripted into acting as an assistant to a machine'': the warehouse worker paced and scored by the system they tend, the driver rated across nine apps. The diagnostic is not the gadget's cleverness. It is who decides where, when, and how it is used. The same device serves one arrangement and consumes the other.
\end{definitionbox}

That axis maps cleanly onto the taxonomy this book set out in Chapter~\ref{ch:jagged-frontier} and Chapter~\ref{ch:three-models}. A Centaur operator keeps the machine subordinate; a Self-Automator risks handing the machine the wheel and becoming its peripheral. Whether an agentic deployment produces one or the other is a question about arrangements rather than silicon, which is exactly Doctorow's second move. He calls the opposing belief inevitabilism, the insistence that there is one way to do things and that any harm it causes is nobody's fault, an echo of Thatcher's ``there is no alternative'' \parencite{Doctorow2023Neoliberalism}. The word predates him; Zuboff had already turned it on tech-determinist rhetoric \parencite{Zuboff2019SurveillanceCapitalism}. His contribution is a worked refutation. Take a cash register: the same automation can gut a clerk's wages or free their afternoon, and which it does turns entirely on the labour arrangement around it. ``The social arrangements of technology are a choice, not an inevitability,'' he writes, and inevitabilism's whole function is to forbid the conversation about that choice, the same conversation the Luddite caricature forbids.

The purest present-day statement of the centaur claim is made from the employer's chair. In ``The Self-Driving Company'' \parencite{MasadKennedy2026}, Amjad Masad and Scott Kennedy report that self-driving ``turns doers into directors'', and put the claim in the workforce's own supposed voice: ``People don't feel like they've been automated. They feel like they've been promoted''. Read charitably, that is a centaur outcome described from the inside: judgement kept, the machine subordinate, the afternoon freed. Read honestly, it is the view from the employer's telemetry, self-reported and unaudited, of AI-fluent engineers at an AI company during a hiring expansion, which is the one population most placed to be promoted rather than displaced. Doctorow's reverse-centaur is about everyone not in that room: the worker paced by a system they did not choose, in a firm that is shedding headcount, not adding it. The post proves the transition can be a promotion under those conditions. It cannot show that it will be, generally.

Doctorow is equally hard on his own side. Borrowing Lee Vinsel's idea of ``criti-hype'', criticism that both feeds and feeds on hype, he warns that repeating an industry's grandest claims, even to attack them, helps sell them \parencite{Vinsel2021}. The warning bites on the numbers, because the numbers are the argument. His account of the bubble runs on a few mechanics. The pitch is aimed at investors rather than users, because a growth story lets a firm mint its own acquisition currency in inflated stock. The accounting circulates: compute credits Microsoft extends to OpenAI get booked by one party as revenue and by the other as investment, the same money counted twice. The market size is fabricated by multiplying exposed headcounts by their salaries and a guessed willingness-to-pay, which is how a bank reaches a sixteen-trillion-dollar figure. The concentration is real and fragile: on Ed Zitron's July 2025 accounting, seven firms carry some thirty-five per cent of the value of all US stocks, and one bad quarter can erase hundreds of billions in a day, as Nvidia's post-DeepSeek fall in 2025 showed. Underneath sits Stein's Law, which Doctorow quotes flat: ``If something cannot go on forever, it will stop'' \parencite{Doctorow2026ReverseCentaur}.

The mechanic that matters for a builder is where the money actually comes from. Consumer novelties do not fund this; the bubble is inflated by a specific promise sold to employers, that the technology will cut the wage bill. The podcast compresses it into a single line about the business model: ``How is it going to increase productivity? By replacing labor'' \parencite{LynskeyDunt2026Luddites}. That is the 1812 dispute restated. And when the replacement misfires, the cost is routed to whoever is cheapest to blame. Doctorow's example is the Air Canada chatbot that invented a bereavement refund policy; the airline disclaimed its own bot, and a tribunal ordered it to pay the fare anyway, establishing that legal liability, rather than the vendor's assurance of trustworthiness, was the mechanism that worked. Beneath many ``AI'' products he finds people: the cashierless shop whose cameras turned out to route to reviewers overseas, the billion-dollar coding tool later described as a staffing agency for hundreds of human programmers. His prescription for creative workers is the WGA's 2023 settlement, under which writers kept the tools and won the right to use them on their own terms, with no forced adoption and no pay cut. In his phrase, ``their bosses tried to turn them into reverse centaurs, and they won the right to be centaurs instead.'' It is the sixpence-per-yard demand restated with a union behind it, a claim on the share of the yield rather than a demand to switch the machine off.

This book has to take his next chapter personally, because it is about what this book builds. His charge sheet against agentic AI is serious and deserves answering without a brochure. He argues that the label is a marketing relabel to re-inflate a tiring bubble; that the base capability is thin, early agents reduced to screenshotting a page and guessing where to click; that the commercial web is deliberately built to be illegible, its owners suing scrapers and killing the apps that would let a buyer compare prices, so an agent working for a user runs against the interests of every site it must cross; and that the proposed fix of agent-legible standards repeats the failed semantic web, because parties with opposed interests will never agree on a shared description and cheating is rewarded. His summary is the sharpest sentence written about the field: the agentic AI industry, he writes, now insists that everyone ``rearrange all of their affairs to the benefit of AI companies, and to their own detriment.'' He points at the hidden humans, the Cruise cars overseen by more than one engineer apiece, and at the fragility, a traffic cone that immobilises a robotaxi.

The honest reply owes nothing to marketing. It is the register this book keeps, and on several of his charges he is simply right. Where the mesh diverges from what he describes, it diverges by construction. Identity is verified through a cryptographic spine rather than inferred from a screenshot, so an agent presents as itself and cannot pass as a person; disclosure of agency is a stated obligation carried by a visible badge, an obligation of the system rather than a courtesy that evaporates the quarter it costs money. Irreversible acts escalate to a human by default rather than proceeding on the machine's confidence. Capability is published in falsifiable tiers instead of asserted. Every action leaves a receipt an adversary can check. Where he catalogues buried human labour, the substrate's reply is to name the humans in the loop and log their interventions, so that a supervisor is a recorded authority with the time and standing to act rather than an unnamed absorber of blame.

On his hardest point this book concedes the ground. Where an agent acts for a user against a counterparty that profits from illegibility, the substrate cannot legislate the counterparty's cooperation, and it should not pretend to; a mesh can make its own agent honest and cannot make a hostile airline's website legible. What it can do is bound the agent's authority and keep its conduct auditable, which is a claim that can be tested rather than a virtue that must be trusted. Doctorow's real target is undisclosed autonomy and buried labour, the two failure modes his book documents again and again. The answer to that is to make disclosure and receipts structural properties of the system, and then to invite the receipts to be checked, which is the discipline the next two sections describe.

\section{What Rhymes, and What Does Not}

The rhyme should be stated without hysteria, because the differences matter as much as the echoes. What rhymes is the shape of the problem. The gains of mechanisation flowed to the mill owners; there is no evidence yet that the beneficiaries of automation this time intend to share the dividend, and no mechanism obliging them to. The channels for raising the grievance are thin. The Combination Acts have no modern statute-book equivalent, and yet a platform's terms of service and a model provider's usage policy occupy the same structural slot: rules written by one party, binding on the other, arriving without appeal and degrading once the users are captured \parencite{Doctorow2025Enshittification}. The Combination Acts forbade the workmen from combining to answer their employers; a usage policy that can be revised overnight forbids very little in law and achieves much the same asymmetry in practice. The information runs one way besides, so a worker cannot calibrate trust in a system they are not allowed to see, any more than a cropper could audit the ledger of the mill that undercut him. And the displacement can outrun the adjustment, which was the whole tragedy at York.

\begin{evidencebox}[title={Early Signals, Read as Preliminary}]
The 2024--26 evidence is real, mixed, and mostly working-paper stage, and this book flags it as such. Aggregate unemployment shows no clear AI-driven break yet, in the OECD's synthesis and in Anthropic's own telemetry-based index \parencite{OECD2023,AnthropicLabor2026}. Underneath the aggregate, a signal is forming at the entry level: \textcite{Brynjolfsson2025Canaries} find a roughly sixteen per cent relative employment decline for workers aged 22--25 in the most exposed occupations since late 2022, concentrated where AI automates rather than augments, and \textcite{Hampole2025} find the exposure falling on high-wage rather than routine work, reversing the older pattern. The task model warns this is a balance rather than a verdict: automation displaces labour from some tasks while new tasks can reinstate it elsewhere, and the net turns on which of the two moves faster \parencite{Acemoglu2019,Autor2015}. That balance is set by institutions rather than by the technology, which is the finding a managed transition can act on.
\end{evidencebox}

What does not rhyme is the thing that makes this chapter an argument rather than a lament. In 1812 an adjustment mechanism had to be legislated: a levy passed, a phasing mandated, a fund administered, each of which Parliament could and did refuse. The substrate this book describes is programmable, so the mechanisms can be built into how the system works rather than petitioned for after the damage. The just-transition literature that grew up around decarbonisation is the anchor for treating the change as governable at all, with its insistence that who bears the cost of a structural shift is a political choice open to design \parencite{ILO2015,NewellMulvaney2013,AutorMindellReynolds2022}. The Luddites had to ask a hostile legislature to build them a transition. A programmable substrate can ship one.

That claim would be cheap if this book exempted its own systems from suspicion, so it does not. The register that scored this ecosystem's code found drift in its own tree this sprint, small and specific. A status indicator was hardcoded to report green regardless of the underlying state. A flag announced that writes had been performed when none had. Capability claims sat a tier above the evidence for them. Twelve times across the sprint's three waves, work reported closed was caught open by a verifier holding a refutation mandate, and sent back (Chapter~\ref{ch:gap-register}). None of this is catastrophic and all of it is the point: the drift toward fabricated status and undisclosed non-action showed up as a measured property of an honestly built system, caught in this book's own tree, and the canaries and receipts are what surfaced it. It is no story about somebody else's deployment. Every one of those defects is a miniature of the exact harm Doctorow catalogues at scale, a machine reporting a state it is not in and a human being told a job is done that is not. The system caught them because it was built to leave receipts; a system built to reassure would have shipped all six green. Unmanaged, that is the direction of travel, and the preview is already running, at small scale, in this book's own repository.

\section{Adjustment Mechanisms, Built In}

Set the two lists side by side and the case for the moment writes itself. The Luddites asked for four things, and every one of them has a substrate-level analogue that this ecosystem is building, at tiers worth stating honestly rather than flattering.

They asked for advance notice and a say, and their ultimatum letters were a crude disclosure mechanism, a warning to reverse course before escalation. The analogue is disclosure of agency: an identity spine that every component treats as a requirement, and agent badges that mark an agent as an agent so it cannot pass as human. These are shipped and integrated. The right to steer, which the croppers never had, appears as the right to interrupt: direct-control surfaces, push-to-talk to a running actor, and in-headset approve or deny, shipped at their stated tiers, with escalation-by-default for irreversible acts carried by an authority-class scheme that presently stands alone and is not yet wired through, which the register records rather than hides (Chapter~\ref{ch:gap-register}).

They asked for a share of the machine's yield, the sixpence on every yard. The analogue is value capture where the work happens: the payment ledger, the pod economics, and Bitcoin-anchored provenance trails that let an adversarial party verify an agent's record against a public chain rather than trust a company's word. The sixpence failed because it depended on a state willing to collect and disburse it, and the state declined. A ledger that settles value to the party that produced it, and anchors the proof where anyone can check it, is the same intent built low enough in the stack that no vote is required to switch it on. They asked not to be thrown on the scrap heap, for alternative employment when the old trade closed. The analogue is a role for the displaced coordination class rather than its redundancy, the judgment broker of Chapter~\ref{ch:role-transformation}, whose work of keeping a mesh coherent is precisely the new task that reinstates human labour above the automated one. The cropper's skill was destroyed with no successor role in view; the coordinator's skill has a successor role named and staffed, which is the difference between a scrap heap and a promotion.

\begin{keyinsight}[title={Same Argument, Societal Scale}]
Chapter~\ref{ch:change-problem} made this case for a single organisation: that automation lands well or badly depending on the disclosure, consent, and value-sharing arrangements built around it, and that these can be designed rather than left to drift. This chapter is that argument at the scale of a labour market. The mechanism is identical; only the blast radius changes.
\end{keyinsight}

The honesty about tiers is the mechanism working out loud. Some of these adjustments are shipped and load-bearing, some stand alone and await integration, and some exist as code that is never yet instantiated, all of it graded in the register rather than asserted in a pitch. That grading is the difference from 1812 stated concretely. The Luddites' transition depended on a vote that went against them, held once and lost. A substrate's transition mechanisms are properties that can be checked continuously, and when one is missing the receipt says so. Doctorow's warning is the design constraint that keeps them honest: any protection that is a mere feature will be quietly removed the first quarter it costs money, which is why disclosure, escalation, and auditability are written as obligations of the system and not as options in a settings menu.

\section{Neither Halt Nor Faith}

The position this chapter argues for is neither of the two it is usually offered. It is not the halt that the caricature-Luddite is accused of wanting, a refusal of the machines that no serious source in this chapter actually demanded. It is not the faith the inevitabilist requires, the surrender of the distributional question to an implacable curve. It is the wager this book has already made in Chapter~\ref{ch:2022-wager}: falsifiable commitments, a two-quarter window, published whichever way they resolve.

The Luddites' real tragedy was not that they lost to the machines. Machines that shear cloth better than a hand can are not a defeat worth mourning. Their tragedy was that no institution existed which could have let them win anything at all: no transition fund, no phased rollout, no share of the yield, no channel to raise the grievance that did not end at a gallows. Building that institution, small and checkable first, is the work of the rest of this book. The honest register keeps it from becoming a promise of virtue. A system that caught its own status light lying green does not get to swear it will never drift. What it can offer is that the next drift will be visible and dated when it comes, which is the whole of the difference between a managed transition and an endured one. Stay committed to progress, the podcast's closing counsel runs, and keep the frame humanistic so the progress does not crush the person it is supposed to serve \parencite{LynskeyDunt2026Luddites}. The frames at Rawfolds were broken by men who had asked for that and been handed a rope. This time the tools to answer them are ours to build, and the receipts to prove we did are already in the tree.
